% =====================================================================================
% sections/discussion.tex: Discussion and Limitations, and Conclusion (written at assembly).
%
% Sources (paths relative to session 4aa2c028's scratchpad); every number below also appears in
% another section or comes from one of these:
%   wt_sp2027/SP2027_PLAN.md   status of C2 (V1: L40S timings, Llama mode C, real-weight bytes,
%                              GPU verifier), F2 (OPT-6.7B and Qwen3-4B smoothing pending), F3
%                              (real-weight end-to-end run), F4 (memory), F6 (L40S run), G3, G4, G5
%   wt_sp2027/docs/improvements/C2_v1_proof_size.md Sec. 8 (real smoothed OPT-125M at 512 tokens:
%                              82.2 -> 48.9 MB in mode C, 1.68x; 1.77x in Kpre; 7 residual adds stay
%                              clear under P3; 2,048-token real-weight cells in the L40S package)
%   wt_sp2027/docs/improvements/F2_quality.md (OPT-6.7B 36.0 against 27.1 with a scalar gain)
%   sections/*.tex             the limitations stated there (Rem. rem:scope, Sec. sec:v1
%                              "Limitations", Sec. sec:eval-reexec, Sec. sec:setup-proof "Scope")
% Labels defined here: sec:discussion, sec:conclusion.
% =====================================================================================

\section{Discussion and Limitations}
\label{sec:discussion}

\para{Measurements.} The base implementation was measured on an NVIDIA L40S and an AMD EPYC 9334,
and each improvement and V1 on an RTX~2080~Ti and a Xeon Silver 4114, against the previous code on
the same queries. The improvements change no byte of the proof, so the proof sizes stand, but the
times of the frozen code on the L40S are not yet measured, and several cells of
Tables~\ref{tab:eval-llm}, \ref{tab:eval-v1} and~\ref{tab:eval-compare} are empty for that reason
\todo{L40S: final run of the frozen code; fill every \texttt{L40S} cell}. The 7B and 13B models were
measured as one to four of their blocks on the 2080~Ti, and full-model figures for them come from
the L40S run or from the byte model. Peak memory is not yet reported \todo{F4}.

\para{Real weights and model quality.} The costs were measured with random int8 weights. They
depend only on the shapes, except for the compressed size of the claims and of V1's values. On
the smoothed OPT-125M at 512 tokens, real activations compress less well than random ones as int8
values, and V1 shrinks the proof by 1.68$\times$ in mode $\sC$ (82.2 to 48.9\,MB) and
1.77$\times$ in $\Kpre$, less than with random weights (1.94--2.06$\times$ for GPT-2 at 512
tokens), also because seven of its residual operations have multipliers above $2^{29}$ and stay
clear under (P3) \todo{real-weight cells at
2{,}048 tokens, and the end-to-end run of Qwen3-4B and OPT-6.7B (F3)}. Smoothing brings OPT-125M and
OPT-1.3B within a factor of 1.03 of fp32. We have not yet run it on OPT-6.7B, whose integer model
with a scalar gain is 33\% above fp32 (36.0 against 27.1), or on Qwen3-4B, because both need more
memory than our 2080~Ti setup provides. Soundness holds for any committed integer model; these
numbers concern only how well that model serves.

\para{V1.} V1 is an option, not a replacement. It halves the proof at 2{,}048 tokens but makes the
prover about ten times slower on the 2080~Ti and the CPU verifier 1.6--1.9$\times$ slower. Llama-2-7B
in mode $\sC$ does not fit on the 11\,GB card with the GKR, so V1 has not been timed in that
configuration, and the full 7B and 13B models have only the byte model's prediction
\todo{L40S: V1 on the full Llama-2-7B, including mode $\sC$, Llama-2-13B, OPT-6.7B and Qwen3-4B}.
Mode $\sK$ and the streaming GPU verifier do not support V1 yet, and V1 assumes one requantisation
multiplier per operation, which excludes the per-row weight scales that would bring OPT-125M to
within $\times1.00$--$1.01$ of fp32.

\para{Proof size and the verifier.} Even with V1, the proof grows with the prompt, because it
carries one value for every output of every weight operation. For Llama-2-7B it reaches the size of
the int8 weights at about 2{,}050 tokens without V1 and, by the byte model, at about 4{,}500 tokens
with it. Only committing the activations and proving the operations without weights, as zkLLM does,
would remove this growth; we leave it to future work. The verifier recomputes attention, so its cost
grows with the square of the prompt length. A CPU client that holds the weights saves time over
re-execution at 64 tokens and ties with fp32 re-execution at 2{,}048 tokens, and on a GPU
re-execution is faster than our verifier. The case for the protocol is therefore a client that does
not hold the weights or the hardware to run them, not a faster way to compute an answer.

\para{Weights and trust.} The protocol is not zero-knowledge: the claims of about $\lceil k/T\rceil$
queries reveal a weight matrix with rows of length $k$, so mode $\sC$ targets open-weight models. For
private weights it shares the limitation that Hollow-LLM points out for zkSNARKs~\cite{hollowllm}:
the commitment binds the provider to one fixed model, but a client that sees neither the weights nor
a range proof on them does not learn which model that is, and weights with a degenerate structure
keep the declared shapes while being cheap to evaluate. With an untrusted commitment, the bound (P4) of V1 and the claim range of v0
must be published with the model, and the time of the setup proof is not yet measured
\todo{L40S: setup prover and verifier time}. The exact sampler is a fixed approximation of softmax
sampling whose distance to it we have not measured, the proof does not show that the client's seed
was random, and nucleus sampling is not implemented.

\para{The attack and the theorem.} The lower bound needs condition (H), which we establish for
random ensembles with independent rows and on a constructed worst-case network, not for a trained
network, whose rows are correlated; on trained networks we report the detection of the best attacks
we found, which bounds each sampler's worst-case detection from above. We do not claim that the worst-case instance embeds in a
transformer, since the LayerNorm before the MLP prevents feeding it the required patterns
(Remark~\ref{rem:scope}). The natural ranges in our range-respecting attacks were computed on test
queries that include the attacked ones, and the adaptive game against value-aware samplers with
held-out ranges is not yet run \todo{G3}. We deliberately do not develop a method for steering the
responses of instruction-tuned assistants.

\section{Conclusion}
\label{sec:conclusion}

A proof of inference that spot-checks an unencoded execution trace is not fully sound. Editing one
value and recomputing the rest honestly turns the forged trace into the honest trace of a model that
differs in one row, and finding that row is a search for a hidden object, in which no sampling rule,
adaptive or value-aware, beats uniform sampling in the worst case. The object that must be encoded is
therefore the model, not the trace. Our defence checks every weight operation of an exact integer
model with Freivalds' algorithm against a Reed--Solomon and Merkle commitment. It needs only a hash
function and a linear code, stays sound against an untrusted committer with a one-time setup proof,
binds the whole computation under Fiat--Shamir, and proves greedy and sampled generation. Its prover
is far faster than published zkSNARK provers. Its cost is a proof that grows with the prompt, which
V1 halves at 2{,}048 tokens but does not remove.
